\chapter{Описание стенда и подготовка}
\label{ch:stand}

\section{Codepen как онлайн-редактор}

\textbf{Codepen} (\texttt{codepen.io}) --- браузерный редактор фрагментов
HTML, CSS и JavaScript, созданный Крисом Койером и Тимом Сэбэтом в
2012~году. Сервис ориентирован на быстрые демонстрации UI-приёмов и
учебные задания: каждый <<pen>> --- это набор из трёх панелей
(\texttt{HTML}, \texttt{CSS}, \texttt{JS}) и панели \texttt{Preview},
которая мгновенно перерисовывается при каждом изменении кода.

В рамках лабораторной работы~12 используется только базовый функционал:
форк готового pen-а, редактирование разметки/стилей и сохранение
результата под собственным аккаунтом.

\section{Регистрация}

На странице \texttt{codepen.io} в правом верхнем углу --- кнопка
\texttt{Sign~Up}. Доступны три способа:
\begin{itemize}
  \item \texttt{Continue~with~GitHub} --- быстрая авторизация через
        GitHub-аккаунт (рекомендуется для разработчиков);
  \item \texttt{Continue~with~Google} --- через Google-аккаунт;
  \item Email~+~пароль --- классическая регистрация с подтверждением
        по почте.
\end{itemize}

Бесплатного тарифа \texttt{Free} достаточно: лимит --- неограниченное
число публичных pen-ов, до трёх частных за всё время существования
аккаунта; кодпен-эссеты (картинки, шрифты) могут быть подключены по
внешним URL.

% — Скриншот: окно регистрации Codepen —
\begin{figure}[h]
  \centering
  \includegraphics[width=0.85\textwidth]{img/09_codepen_save}
  \caption{Интерфейс редактора Codepen: панели HTML/CSS/JS и Preview}
  \label{fig:codepen_ui}
\end{figure}

\section{Структура pen-а}

Каждый pen состоит из четырёх взаимосвязанных панелей:
\begin{itemize}
  \item \texttt{HTML} --- разметка тела документа. Codepen
        автоматически оборачивает её в \texttt{<!DOCTYPE~html>},
        \texttt{<html>}, \texttt{<head>}, \texttt{<body>}: писать
        обвязку вручную не требуется.
  \item \texttt{CSS} --- стили; по умолчанию подключаются как
        внутренний \texttt{<style>}.
  \item \texttt{JS} --- JavaScript-код, исполняемый после загрузки
        DOM (в текущей лабораторной не используется).
  \item \texttt{Preview} --- iframe с отрисованным результатом;
        обновляется автоматически или по \texttt{Ctrl+S}.
\end{itemize}

Через кнопку \texttt{Settings} к pen-у подключаются внешние ресурсы:
\begin{itemize}
  \item \texttt{HTML~\(\rightarrow\)~Stuff~for~<head>} --- произвольные
        теги в \texttt{<head>} (например, \texttt{<meta>}, \texttt{<link>}
        на шрифты);
  \item \texttt{CSS~\(\rightarrow\)~Add~External~Stylesheets} --- список
        URL внешних таблиц стилей; в стартовых pen-ах текущей
        работы здесь уже подключены Material~Design~Lite и базовая
        таблица Netology, поэтому в собственный CSS их повторять
        не нужно.
  \item \texttt{CSS~\(\rightarrow\)~CSS~Preprocessor} --- выбор препроцессора
        (\texttt{None}, \texttt{SCSS}, \texttt{LESS}, \texttt{Stylus});
        в этой работе остаётся значение \texttt{None}.
\end{itemize}

\section{Fork-workflow}

Стартовые pen-ы методички принадлежат автору \texttt{Andrei-Andreev-the-builder}
и доступны студенту только в режиме просмотра. Чтобы внести правки,
нужно создать собственную копию --- \emph{форк}:

\begin{enumerate}
  \item Перейти на стартовый pen по ссылке \texttt{clck.ru/...} из
        методички (таблица соответствий приведена в разделе~\ref{ch:practice}).
  \item Войти в аккаунт Codepen.
  \item Нажать кнопку \texttt{Fork} в правом верхнем углу страницы
        pen-а. Codepen создаст новый pen с тем же содержимым, но
        под текущим пользователем, и автоматически откроет его в
        режиме редактирования.
  \item Внести правки в панели \texttt{HTML} (для блоков~1.1, 1.2)
        и/или \texttt{CSS} (для блоков~2.1, 2.2). Preview
        обновляется автоматически.
  \item Нажать \texttt{Save} (или \texttt{Ctrl+S}) --- pen
        получит уникальный URL, который остаётся за студентом
        навсегда. Этот URL вставляется в раздел \texttt{<<Готовые
        форки в Codepen>>} главы~\ref{ch:practice}.
\end{enumerate}

% — Скриншот: кнопка Fork → Save в Codepen —
\begin{figure}[h]
  \centering
  \includegraphics[width=0.85\textwidth]{img/09_codepen_save}
  \caption{Кнопки \texttt{Fork} и \texttt{Save} в правой верхней панели Codepen}
  \label{fig:codepen_fork}
\end{figure}

\section{Параллельная локальная разработка}

Codepen удобен для коротких демо, но плохо подходит для редактирования
вне браузера или для контроля версий через git. Поэтому в каталоге
\texttt{lab12/codepen-solutions/} хранятся локальные копии каждого
из четырёх решений: четыре подкаталога (\texttt{1.1-html-tags/},
\texttt{1.2-lists-tables/}, \texttt{2.1-selectors/},
\texttt{2.2-text-styling/}), каждый с собственным \texttt{index.html}
и --- для блоков, требующих CSS, --- \texttt{styles.css}.

Все четыре изображения, требуемые методичкой (три фотографии гейзеров
Камчатки для блока~1.1 и фон цитаты для блока~2.2), скачаны из
методички Яндекс.Диска и положены в \texttt{codepen-solutions/img/};
ссылки в HTML --- относительные (\texttt{../img/<имя>.jpg}).
Это позволяет открывать локальные страницы офлайн: для рендеринга
требуется только CDN с шрифтами Google~Fonts, Material~Design~Lite
и общими таблицами стилей Netology (\texttt{netology-code.github.io}).

Запуск любого решения:
\begin{lstlisting}[language=bash]
xdg-open lab12/codepen-solutions/1.1-html-tags/index.html
xdg-open lab12/codepen-solutions/1.2-lists-tables/index.html
xdg-open lab12/codepen-solutions/2.1-selectors/index.html
xdg-open lab12/codepen-solutions/2.2-text-styling/index.html
\end{lstlisting}

После проверки локального рендера правки переносятся в
соответствующий форк Codepen.

\endinput
